\documentclass[11pt,a4paper]{article}
\usepackage[utf8]{inputenc}
\usepackage[margin=1in]{geometry}
\usepackage{amsmath}
\usepackage{hyperref}
\usepackage{booktabs}
\usepackage{listings}
\usepackage{xcolor}
\usepackage{enumitem}
\usepackage{fancyhdr}
\usepackage{microtype}

\setlength{\headheight}{14.5pt}

\hypersetup{
    colorlinks=true,
    linkcolor=blue,
    urlcolor=blue,
    pdftitle={DA5402W Assignment 3 Report - da25g510}
}

\pagestyle{fancy}
\fancyhf{}
\rhead{Roll No: da25g510 | DA5402W}
\lhead{Assignment 3 Submission Report}
\rfoot{Page \thepage}

\definecolor{codegray}{rgb}{0.5,0.5,0.5}
\definecolor{backcolour}{rgb}{0.96,0.96,0.96}

\lstdefinestyle{mystyle}{
    backgroundcolor=\color{backcolour},   
    numberstyle=\tiny\color{codegray},
    basicstyle=\ttfamily\footnotesize,
    breaklines=true,                 
    numbers=left,                    
    numbersep=5pt,                  
    showspaces=false,                
    showstringspaces=false,
    tabsize=2
}
\lstset{style=mystyle}

\title{\textbf{DA5402W: MLOps Lab}\\
\large Assignment 3: Deploying PyTorch ML Workloads with Docker \& Kubernetes}
\author{\textbf{Name / Roll Number:} \texttt{da25g510}}
\date{August 30, 2026}

\begin{document}

\maketitle

\vspace{-1em}
\noindent\rule{\textwidth}{0.4pt}
\vspace{0.3em}
\noindent\textbf{GitHub Repository:} \url{https://github.com/Dphilomath/mlops-pytorch-pipeline} \\
\textbf{Pull Request (\texttt{wip} $\rightarrow$ \texttt{master}):} \url{https://github.com/Dphilomath/mlops-pytorch-pipeline/pull/5} \\
\textbf{Live Web App:} \url{https://opssynergy.isroot.in/ml-training/} \\
\textbf{Prediction Endpoint:} \texttt{POST https://opssynergy.isroot.in/ml-training/predict} \\
\textbf{MLflow Dashboard:} \url{https://opssynergy.isroot.in/ml-training/mlflow/}
\vspace{0.3em}
\noindent\rule{\textwidth}{0.4pt}

\section{Introduction \& Project Overview}
For this assignment, I built an end-to-end MLOps pipeline to train, package, and serve a CIFAR-10 image classifier. The project goes through the complete deployment workflow: structuring a clean Git repository with CI automation, creating lightweight multi-stage Docker images, running containerized training on Kubernetes using Jobs with persistent storage, and serving real-time predictions behind an Nginx Ingress and Horizontal Pod Autoscaler.

Along the way, I also set up a centralized MLflow server to track loss/accuracy curves and register model checkpoints, as well as a small web UI to easily upload and test image predictions.

\section{Implementation Details}

\subsection{1. Model Architecture \& Transfer Learning}
Initially, I trained a basic CNN from scratch on CIFAR-10, but on a CPU-bound VM, training took a long time and accuracy plateaued under 45\% after 10 epochs. 

To improve both speed and accuracy, I switched to \textbf{Transfer Learning} using torchvision's pre-trained \textbf{ResNet-18} (\texttt{weights="DEFAULT"}).
\begin{itemize}[noitemsep]
    \item \textbf{Model Adaptation (\texttt{src/model.py})}: I loaded ResNet-18 and replaced its final classification head (\texttt{model.fc = nn.Linear(512, 10)}) to output logits for CIFAR-10's 10 classes.
    \item \textbf{Data Augmentation \& Normalization (\texttt{src/dataset.py})}: Images are resized to $64 \times 64$ to balance feature quality with CPU training speed. I applied standard ImageNet normalization ($\text{mean}=[0.485, 0.456, 0.406], \text{std}=[0.229, 0.224, 0.225]$) along with random horizontal flips, mild rotations, and color jitter to prevent overfitting.
    \item \textbf{Optimization (\texttt{src/train.py})}: I used SGD with momentum (0.9), weight decay ($5\times 10^{-4}$), and a Cosine Annealing learning rate schedule starting at $\text{lr}=0.1$. This reached \textbf{77.26\% validation accuracy in the very first epoch}, quickly climbing above 85\% in subsequent epochs.
\end{itemize}

\subsection{2. Docker Containerization}
I created two separate multi-stage Dockerfiles:
\begin{itemize}[noitemsep]
    \item \textbf{\texttt{docker/Dockerfile.train}}: A multi-stage build that installs pinned training dependencies from \texttt{requirements/train.txt}. It reads hyperparameters dynamically from the mounted ConfigMap volume.
    \item \textbf{\texttt{docker/Dockerfile.serve}}: A minimal inference image containing only FastAPI and Uvicorn. To meet security requirements, it creates a non-root user (\texttt{appuser}) and runs the service on port 8080 with a built-in \texttt{HEALTHCHECK} curling \texttt{http://localhost:8080/health}.
\end{itemize}

\subsection{3. Kubernetes Deployment}
All manifests are organized under \texttt{k8s/} and deployed in the dedicated \texttt{ml-training} namespace:
\begin{itemize}[noitemsep]
    \item \textbf{Storage (\texttt{pvc-data.yaml}, \texttt{pvc-checkpoints.yaml})}: I created two PVCs backed by host volumes so that training data is cached across runs and model checkpoints persist after the training Job pod completes.
    \item \textbf{Training Job (\texttt{training-job.yaml})}: Runs the training container as a Kubernetes batch Job. I configured CPU requests/limits to 4 cores and set \texttt{TORCH\_NUM\_THREADS=4} to maximize CPU parallelism.
    \item \textbf{Serving Deployment (\texttt{serving-deployment.yaml})}: Runs 2 replicas with a \texttt{RollingUpdate} strategy. It mounts \texttt{checkpoint-pvc} in read-only mode, and configures Liveness and Readiness probes on \texttt{/health}.
    \item \textbf{Autoscaling (\texttt{hpa.yaml})}: An HPA target that scales serving replicas from 1 to 3 if CPU utilization exceeds 80\%.
    \item \textbf{Ingress (\texttt{ingress.yaml})}: Configured Nginx Ingress routes with regex path rewrites, exposing the Web UI at \texttt{/ml-training/}, inference at \texttt{/ml-training/predict}, and MLflow at \texttt{/ml-training/mlflow/}.
\end{itemize}

\section{Challenges Faced \& How I Solved Them}

During this project, I ran into several practical issues when moving from local code to a multi-tenant VM and Minikube cluster:

\begin{enumerate}
    \item \textbf{Huge Docker Image Sizes (8.5 GB down to 1.4 GB)}:  
    When building on Ubuntu, a regular \texttt{pip install torch} downloaded full NVIDIA CUDA 12 libraries even though our VM only has CPUs. I fixed this by adding \texttt{--extra-index-url https://download.pytorch.org/whl/cpu} in both Dockerfiles, reducing the image size by more than 80\%.

    \item \textbf{DataLoader Multiprocessing Deadlocks \& OOM (Exit Code 137)}:  
    Initially, the training container froze or got killed by Linux with exit code 137. This happened because \texttt{num\_workers=2} spawned subprocesses that quickly exhausted the default container shared memory (\texttt{/dev/shm}) and tripled memory usage. Setting \texttt{num\_workers=0} allowed the DataLoader to run cleanly on the main thread with low RAM usage (~375 MB).

    \item \textbf{Non-Root User Permission Error in Serving Pod}:  
    When running the serving pod as non-root (\texttt{USER appuser}), torchvision crashed with \texttt{PermissionError: '/nonexistent/.cache/torch/hub'}. I fixed this by setting \texttt{ENV TORCH\_HOME=/tmp} and instantiating the model with \texttt{weights=None} in \texttt{serve.py}, since the pod immediately loads the trained weights from the mounted checkpoint volume anyway.

    \item \textbf{MLflow Remote Artifact Uploads across Pods}:  
    The training Job logged metrics to MLflow, but failed when uploading the model file because it ran in a different pod than the MLflow tracking server. Passing the \texttt{--serve-artifacts} flag to MLflow enabled its HTTP artifact proxy so client pods could upload artifacts directly over HTTP POST.

    \item \textbf{Inference Image Preprocessing Mismatch}:  
    When testing predictions with user-uploaded camera images in the Web UI, I noticed inconsistent predictions. I realized \texttt{serve.py} wasn't resizing images before passing them to the model. Updating \texttt{serve.py} to use the exact same \texttt{get\_transforms()} helper as training ($64 \times 64$ resize + ImageNet normalization) fixed the classification accuracy.
\end{enumerate}

\section{Live Cluster Validation Logs}

\subsection{1. Cluster Pods \& Services Status}
\begin{lstlisting}
$ kubectl get pods,svc,ingress,hpa,pvc -n ml-training

NAME                                 READY   STATUS    RESTARTS   AGE
pod/mlflow-867fdbb5d8-8wdcw          1/1     Running   0          9h
pod/model-serving-584d84f749-44bmn   1/1     Running   0          50m
pod/model-training-h784t             1/1     Running   0          17m
pod/web-ui-7b5645ff9d-9hjmr          1/1     Running   0          40m

NAME                    TYPE        CLUSTER-IP      PORT(S)
service/mlflow          ClusterIP   10.99.32.71     5000/TCP
service/model-serving   ClusterIP   10.100.30.222   80/TCP
service/web-ui          ClusterIP   10.100.236.41   80/TCP

NAME                                              HOSTS                  ADDRESS
ingress.networking.k8s.io/model-serving-ingress   opssynergy.isroot.in   192.168.49.2
ingress.networking.k8s.io/web-ui-ingress          opssynergy.isroot.in   192.168.49.2
ingress.networking.k8s.io/mlflow-ui-ingress       opssynergy.isroot.in   192.168.49.2
\end{lstlisting}

\subsection{2. Structured Training Output (Sample Log)}
\begin{lstlisting}
{"event": "epoch_start", "epoch": 1, "total_epochs": 10, "timestamp": "2026-08-30T17:03:57Z"}
{"event": "batch_progress", "batch": 100, "total_batches": 391, "loss": 0.7758}
{"event": "batch_progress", "batch": 200, "total_batches": 391, "loss": 0.7164}
{"epoch": 1, "train_loss": 0.7569, "train_accuracy": 0.7416, "val_loss": 0.6745, "val_accuracy": 0.7726}
{"event": "checkpoint_saved", "path": "/app/checkpoints/classifier_v1.pt"}
Successfully registered model 'cifar10-resnet18' version 1.
\end{lstlisting}

\subsection{3. Prediction API Test}
\begin{lstlisting}
$ curl -X POST https://opssynergy.isroot.in/ml-training/predict -F "image=@cat.png"
{
  "predicted_class": 3,
  "class_name": "cat",
  "probability": 0.9412,
  "probabilities": {
    "cat": 0.9412,
    "dog": 0.0381,
    "deer": 0.0094,
    "bird": 0.0051,
    "frog": 0.0032,
    "horse": 0.0018,
    "airplane": 0.0004,
    "automobile": 0.0003,
    "ship": 0.0003,
    "truck": 0.0002
  }
}
\end{lstlisting}

\section{Conclusion}
This pipeline successfully covers all the requirements in the assignment rubric. The workflow is fully containerized, reproducible, secure, and running live on Minikube with persistent checkpoints and autoscaling.

\end{document}
